\documentclass[11pt]{article}
\usepackage[a4paper,margin=1in]{geometry}
\usepackage{amsmath,amssymb,amsthm}
\usepackage{bm}
\usepackage{hyperref}
\usepackage{xcolor}
\hypersetup{colorlinks=true,linkcolor=blue,citecolor=blue,urlcolor=blue}

\title{Hybrid--Smooth Backward Mode for \texttt{TinyMultiHeadLut}\\
\large Mathematical Specification}
\author{LUTorch project notes}
\date{2026--05--29}

\newcommand{\NAP}{N_{\!\text{a}}}
\newcommand{\nalt}{k}
\newcommand{\bts}{\mathrm{ts}}
\newcommand{\Tsoft}{T_{\!\text{soft}}}
\newcommand{\Tsel}{T_{\!\text{sel}}}

\begin{document}
\maketitle

\section{Setup}

We process one slot (one head, one table) at a time; the same math repeats
$H \times \text{tph}$ times in parallel. Inputs:
\begin{itemize}
  \item Input vector $\bm{x} \in \mathbb{R}^{d_{\text{in}}}$.
  \item Anchor pairs $(a_i, b_i) \in \{0, \dots, d_{\text{in}}{-}1\}^2$ for
        $i = 0, \dots, \NAP{-}1$; chosen at module init, fixed thereafter.
  \item Weight matrix $W \in \mathbb{R}^{2^\NAP \times d_{\text{out}}}$, with
        rows $W_r$ indexed by binary patterns $r \in \{0,1\}^\NAP$.
  \item Two learnable scalar temperatures: $\Tsoft, \Tsel > 0$
        (stored as $\log T$; we use $\Tsoft = \exp \log T_{\!\text{soft}}$, etc.).
\end{itemize}

\paragraph{Anchor differences.} For each pair $i$,
\begin{equation}
  d_i \;=\; x_{a_i} - x_{b_i}, \qquad i = 0, \dots, \NAP{-}1.
\end{equation}

\paragraph{Hard ("main") row index.} The bit pattern from sign-of-difference,
MSB-first:
\begin{equation}
  \text{main}(\bm{x}) \;=\; \sum_{i=0}^{\NAP-1} \mathbf{1}[d_i > 0] \cdot 2^{\NAP-1-i}
  \;\in\; \{0, \dots, 2^\NAP-1\}.
\end{equation}

\section{Per--anchor soft score}

We define a smooth surrogate for the sign:
\begin{equation}
  p_i \;=\; \operatorname{sign}(d_i) \cdot \frac{|d_i|}{\Tsoft + |d_i|},
  \qquad p_i \in (-1, 1).
\end{equation}
The absolute value
\begin{equation}
  |p_i| \;=\; \frac{|d_i|}{\Tsoft + |d_i|}
\end{equation}
is a confidence in the sign assignment: $|p_i| \to 0$ for $|d_i| \ll \Tsoft$
(uncertain), $|p_i| \to 1$ for $|d_i| \gg \Tsoft$ (very confident).
$\Tsoft$ sets the scale at which a difference becomes "decisive".

\paragraph{Full row score.} Define a score over rows by summing per-bit
soft contributions:
\begin{equation}
  \bts(r) \;=\; \sum_{i=0}^{\NAP-1} s_i(r), \qquad
  s_i(r) \;=\; \begin{cases}
    +|p_i| & \text{if bit } i \text{ of } r \text{ agrees with } d_i\text{'s sign}, \\
    -|p_i| & \text{otherwise.}
  \end{cases}
\end{equation}
By construction $\bts(\text{main}) = \sum_i |p_i|$ is the maximum over all
rows, and flipping bit $i$ of \emph{main} reduces the score by $2|p_i|$.

\section{Hamming--1 ball restriction}

The full softmax over $2^\NAP$ rows is wasteful (and impractical for backprop)
because the dominant probability mass lives near \emph{main}. We restrict the
softmax to the \emph{Hamming--1 ball} around \emph{main}: the row \emph{main}
itself plus the $\NAP$ rows reachable by flipping a single bit.

\paragraph{Top--$\nalt$ alternatives.}
We further restrict to the $\nalt \le \NAP$ \emph{least confident} bit
positions, i.e.\ the positions with the smallest $|p_i|$:
\begin{equation}
  \mathcal{T} \;=\; \operatorname*{arg\,sort}^{\min}_{i \in \{0,\ldots,\NAP-1\}} |p_i|\;,
  \qquad
  \text{top--}\nalt = \mathcal{T}[0:\nalt].
\end{equation}
For each $j \in \text{top--}\nalt$, the corresponding alternative row is
$\text{main} \oplus 2^{\NAP-1-j}$ (XOR with the bit's place-value).

\paragraph{Score gap.} For each alternative we have
\begin{equation}
  \Delta \bts_j \;\equiv\; \bts(\text{main}) - \bts(\text{alt}_j) \;=\; 2\,|p_j|,
  \qquad j \in \text{top--}\nalt.
\end{equation}
Small $|p_j|$ $\Rightarrow$ small gap $\Rightarrow$ alternative competes.

\section{Forward pass}

\subsection{Logits and probabilities}

Subtracting a constant from all logits doesn't change the softmax, so we
center on $\bts(\text{main})$:
\begin{equation}
  \ell_0 = 0, \qquad
  \ell_{k+1} = -\frac{\Delta \bts_{j_k}}{\Tsel} = -\frac{2|p_{j_k}|}{\Tsel},
  \quad k = 0, \dots, \nalt{-}1.
\end{equation}
The $(\nalt+1)$--way selection probabilities are
\begin{equation}
  \pi_m \;=\; \frac{\exp(\ell_m)}{\sum_{m'=0}^{\nalt} \exp(\ell_{m'})},
  \quad m = 0, \dots, \nalt.
\end{equation}
$\Tsel$ controls how sharp the selection is: $\Tsel \to 0$ collapses to the
hard arg--max (always \emph{main}); $\Tsel \to \infty$ averages uniformly over
the ball.

\subsection{Output}

Letting $r_m$ denote the row indices $(\text{main}, \text{alt}_{j_0}, \ldots,
\text{alt}_{j_{\nalt-1}})$, the per--table output is
\begin{equation}
  \bm{y} \;=\; \sum_{m=0}^{\nalt} \pi_m \, W_{r_m} \;\in\; \mathbb{R}^{d_{\text{out}}}.
\end{equation}
The module sums $\bm{y}$ across the table--per--head dimension to produce
$\bm{o} \in \mathbb{R}^{H \times d_{\text{out}}}$.

\subsection{Special case: $\nalt = 1$}

With a single alternative ($\nalt = 1$ at the smallest $|p|$, call it
$|p_\star|$), the softmax becomes a sigmoid:
\begin{equation}
  \pi_1 \;=\; \sigma\!\left(-\frac{2 |p_\star|}{\Tsel}\right) \;=\;
  \sigma\!\left(-\frac{2}{\Tsel}\,\frac{|d_\star|}{\Tsoft + |d_\star|}\right),
  \qquad \pi_0 = 1 - \pi_1.
\end{equation}
This is the closed--form (Bradley--Terry) top--2 softmax: it returns the same
probability mass on \emph{main} vs.\ its top competitor as a full
$2^\NAP$--way softmax \emph{would}, under the assumption that the third--
through $2^\NAP$--th rows contribute negligibly. Both temperatures enter
multiplicatively in the argument, which is why it's important they appear
in the order shown (a 2026--05--12 correction over an earlier draft that
mis--placed $\Tsel$).

\section{Backward pass (autograd)}

The forward is written entirely in differentiable PyTorch operations modulo
the $\bm{1}[d_i > 0]$ thresholds and the $\arg$-selection that picks the
top--$\nalt$ positions. These two are constant w.r.t.\ infinitesimal changes
in $\bm{x}$ (the sign of $d_i$ and the index of the smallest $|p_i|$ do not
flip), so the gradient is well--defined almost everywhere, and we let
\texttt{torch.autograd} chain it.

\subsection{Weight gradient}

By inspection of $\bm{y}$:
\begin{equation}
  \frac{\partial \mathcal{L}}{\partial W_{r}}
  \;=\; \sum_{m=0}^{\nalt} \mathbf{1}[r_m = r]\, \pi_m \,
  \frac{\partial \mathcal{L}}{\partial \bm{y}}.
\end{equation}
Only the $(\nalt + 1)$ rows in the chosen ball receive gradient; all other
rows get exactly zero. This is the central efficiency win of hybrid--smooth
backward: weight updates are $(\nalt + 1)$--sparse per token regardless of
table size.

\subsection{Input gradient}

The chain is:
\begin{align}
  \frac{\partial \mathcal{L}}{\partial \pi_m}
  &= \frac{\partial \mathcal{L}}{\partial \bm{y}} \cdot W_{r_m}, \\
  \frac{\partial \mathcal{L}}{\partial \ell_m}
  &= \pi_m \left( \frac{\partial \mathcal{L}}{\partial \pi_m}
                 - \sum_{m'} \pi_{m'} \frac{\partial \mathcal{L}}{\partial \pi_{m'}} \right)
     \quad \text{(softmax Jacobian)}, \\
  \frac{\partial \mathcal{L}}{\partial \Delta \bts_{j_k}}
  &= -\frac{1}{\Tsel} \frac{\partial \mathcal{L}}{\partial \ell_{k+1}},
  \qquad
  \frac{\partial \mathcal{L}}{\partial |p_{j_k}|}
  = 2 \frac{\partial \mathcal{L}}{\partial \Delta \bts_{j_k}}, \\
  \frac{\partial |p_{j_k}|}{\partial |d_{j_k}|}
  &= \frac{\Tsoft}{(\Tsoft + |d_{j_k}|)^2}, \\
  \frac{\partial |d_{j_k}|}{\partial d_{j_k}}
  &= \operatorname{sign}(d_{j_k}), \\
  \frac{\partial d_{j_k}}{\partial x_q}
  &= \mathbf{1}[a_{j_k} = q] - \mathbf{1}[b_{j_k} = q].
\end{align}
Only the bit positions in $\text{top--}\nalt$ contribute to
$\partial \mathcal{L}/\partial \bm{x}$; positions outside the ball get
gradient zero. This is intentional: those bits are highly confident
(large $|p_i|$), so flipping them is so unlikely that the loss is
effectively flat in those directions.

\subsection{Temperature gradients}

\begin{align}
  \frac{\partial \mathcal{L}}{\partial \Tsel}
  &= \sum_k \frac{\partial \mathcal{L}}{\partial \ell_{k+1}} \cdot
     \frac{\Delta \bts_{j_k}}{\Tsel^2}, \\
  \frac{\partial \mathcal{L}}{\partial \Tsoft}
  &= -\sum_k \frac{\partial \mathcal{L}}{\partial |p_{j_k}|} \cdot
     \frac{|d_{j_k}|}{(\Tsoft + |d_{j_k}|)^2}.
\end{align}
We parametrize $\Tsoft, \Tsel = e^{\log T}$ and let autograd push the
$\partial T / \partial \log T = T$ factor onto the $\log$-T parameters.

\section{Algorithmic notes}

\subsection{Sequential argmin in lieu of $\operatorname{topk}$}

For small $\nalt$ (typically $1, 2, 3$), we replace
$\operatorname{topk}(|p|, k=\nalt, \text{largest=False})$ with a loop of
$\arg\min{}$s, masking the chosen position to $+\infty$ after each pick:
\begin{equation}
  \begin{aligned}
  &\text{for } k = 0, \dots, \nalt{-}1: \\
  &\quad j_k = \arg\min_i |p_i|, \\
  &\quad |p_{j_k}| \leftarrow +\infty.
  \end{aligned}
\end{equation}
Under \texttt{torch.compile}, this fuses with the surrounding kernel; PyTorch's
\texttt{topk} kernel does not. Empirically at $\NAP{=}6$, $\nalt{=}3$,
$B{=}4096$, $\text{tph}{=}1024$: forward $1.7$\,ms (seq--argmin + compile) vs.\
$21.9$\,ms (topk, eager) --- a $13\times$ kernel--level speedup. Memory drops
$2.7\times$.

\subsection{Why $|p|$ instead of $|d|$ for top selection}

The score gap $\Delta\bts = 2|p|$ is monotonic in $|d|$ (since $|p|$ is monotonic
in $|d|$), so sorting by $|p|$ and sorting by $|d|$ give the same ordering.
Sorting by $|p|$ keeps everything inside the same already--computed tensor,
saving one extra pass.

\section{Numerical considerations}

\begin{itemize}
  \item \textbf{Sign ties.} If any $d_i = 0$ exactly, the "main" bit is taken
        as $0$ (because $\mathbf{1}[d > 0]$ is false). This matches the
        convention used by the standard \texttt{MultiHeadLut} forward; the
        gradient is undefined at the tie point but a.s.\ defined for any
        random initialization.
  \item \textbf{Index--tie selection.} If two positions share the smallest
        $|p_i|$ exactly, $\arg\min$ picks the lower--indexed one
        deterministically (PyTorch convention); \texttt{topk} returns the same
        set but may differ in order. The softmax over the ball is
        permutation--invariant, so the forward output is unaffected.
  \item \textbf{Compile + gradcheck.} \texttt{torch.compile} enables "donated
        buffers" which forbid multiple backward passes. \texttt{gradcheck} in
        tests is wrapped in
        \verb|torch._functorch.config.patch(donated_buffer=False)|.
\end{itemize}

\section{Tests}

\texttt{src/spiky/lutorch/tests/test\_hybrid\_smooth\_nalt.py} verifies:
\begin{enumerate}
  \item \textbf{Forward matches reference.} An independent reference forward
        computes the explicit top--$(\nalt{+}1)$ softmax and matches the
        compiled production forward to \texttt{atol=1e-10} (fp64),
        $\nalt \in \{1,2,3,4\}$.
  \item \textbf{Sequential argmin $=$ topk as sets.} The chosen position
        sets agree up to permutation, $\nalt \in \{1,2,3\}$.
  \item \textbf{Gradcheck.} \texttt{torch.autograd.gradcheck} on
        $(\bm{x}, W, \log T_{\!\text{soft}}, \log T_{\!\text{sel}})$ passes in
        fp64 for $\nalt \in \{1,2,3,4\}$.
  \item \textbf{Eager $=$ compiled.} Output and all gradients agree to
        fp32 noise for $\nalt \in \{1,2,3,4\}$.
  \item \textbf{$\nalt = \NAP$ = full ball.} The specialized
        no--topk branch matches the explicit full--ball reference.
\end{enumerate}

\end{document}
